\documentclass[aspectratio=169,11pt]{beamer}
\usetheme{metropolis}
\usepackage[T1]{fontenc}
\usepackage{lmodern}
\usepackage{hyperref}
\title{Curriculum Vitae}
\subtitle{Academic Profile and Selected Contributions}
\author{Rudi Adha Prihandoko, S.Si., M.Si.}
\institute{Department of Mathematics\\Faculty of Mathematics and Natural Sciences\\Universitas Gadjah Mada}
\date{November 2026}
\begin{document}
\begin{frame}[plain]\titlepage\end{frame}

\begin{frame}{Academic Profile and Education}
\small
\textbf{Academic appointment}\par
Lecturer, Department of Mathematics, Universitas Gadjah Mada, Yogyakarta, Indonesia.
\vspace{0.65em}
\textbf{Education}
\begin{itemize}
\item \textbf{Doctoral studies in Mathematical Sciences (Numerical Mathematics)} --- Australian National University, Australia. \emph{Completion status to be confirmed.}
\item \textbf{M.Si., Mathematics} --- Institut Teknologi Bandung, Indonesia. Thesis: \emph{Kajian tentang Kode Geometri Aljabar sebagai Perluasan dari Kode Reed--Solomon} (2015).
\item \textbf{S.Si.} --- Institut Teknologi Bandung, Indonesia. Thesis: \emph{Kajian tentang Bilangan $p$-adic} (2012)
\end{itemize}
\textbf{Research interests:} numerical analysis; hyperbolic PDEs; shallow-water and sediment transport models; stable numerical methods; coding theory.
\end{frame}

\begin{frame}{Journal Publication and Research}
\small
\textbf{Peer-reviewed publication}\par
Prihandoko, R., Duru, K., Roberts, S., and Zoppou, C. (2024). \emph{On well-posed boundary conditions and energy stable finite-volume method for the linear shallow water wave equation}. \textbf{The ANZIAM Journal}, 66, 181--200.\\
\href{https://doi.org/10.1017/S1446181124000191}{doi:10.1017/S1446181124000191}
\vspace{0.7em}
\textbf{Selected presentations and ongoing research}
\begin{itemize}
\item \textbf{CTAC 2022:} presentation on summation-by-parts and finite-volume analysis of linearised shallow-water equations.
\item \textbf{Manuscript:} \emph{Dual pairing summation-by-parts operator with Godunov flux splitting for the sediment transport model}. \emph{Status to be confirmed.}
\item \textbf{Ongoing:} hyperbolicity and well-posedness of two-dimensional Grass-type sediment transport models.
\end{itemize}
\end{frame}

\begin{frame}{Books, Teaching Modules and Editorial Work}
\small
\begin{enumerate}
\item \textbf{Tutorial Book: Pengantar Analisis Numerik} (2018) --- introductory numerical analysis.
\item \textbf{Tutorial Book: Matematika Komputasi} (2018) --- computational mathematics.
\item \textbf{Geometri Transformasi} (2016) --- transformation geometry.
\item \textbf{Monthly Magazine -- Olimpiade.org} --- mathematics and Olympiad magazine.
\item \textbf{Buku 10 Tahun Olimpiade.org} (2015) --- commemorative publication.
\end{enumerate}
\vspace{0.7em}
\footnotesize Titles and dates supplied by the author. Publication metadata and specific roles (author, editor, or contributor) remain to be verified.
\end{frame}

\begin{frame}{Mathematics Olympiad and Outreach}
\begin{itemize}
\item \textbf{2018:} Jury member, National Science Olympiad (OSN), Mathematics, as listed in the official OSN guide.
\item \textbf{2018--2019:} Involved in provincial Mathematics Olympiad problem preparation and training activities.
\item \textbf{2018:} Mathematics and science teacher development activities in Gunungkidul, Yogyakarta.
\item \textbf{2020:} Participated in mathematics teacher training in Sleman, including higher-order thinking and problem-solving strategies.
\item \textbf{2020:} Contributed to a community-service discussion group on mathematical game application programming.
\end{itemize}
\end{frame}

\begin{frame}{Selected Research and Academic Service}
\begin{itemize}
\item \textbf{2020:} Member of the UGM research project \emph{Teori Bilangan} (Number Theory).
\item \textbf{2017:} Participated in scientific writing training for high-school mathematics and science teachers in Temanggung.
\item \textbf{2022:} Conference research presentation at CTAC, Queensland University of Technology.
\item \textbf{2024:} Publication of research on well-posed boundary conditions and energy-stable finite-volume methods in \emph{The ANZIAM Journal}.
\end{itemize}
\end{frame}

\begin{frame}{Public Academic Sources}
\scriptsize
\begin{itemize}
\item UGM faculty profile: \url{https://mipa.ugm.ac.id/id/2023/10/07/rudi-adha-prihandoko-s-si-m-si/}
\item UGM academic profile: \url{https://prisma.simaster.ugm.ac.id/Rudi}
\item ANZIAM journal article: \url{https://doi.org/10.1017/S1446181124000191}
\item ITB digital library thesis: \url{https://digilib.itb.ac.id/gdl/view_data/kajian-tentang-kode-geometri-aljabar-sebagai-perluasan-dari-kode-reed-solomon/}
\item CTAC 2022 programme: \url{https://cms.qut.edu.au/__data/assets/pdf_file/0019/1231138/ctac-2022-program.pdf}
\item OSN 2018 official guide: \url{https://repositori.kemendikdasmen.go.id/10178/1/Panduan\%20OSN\%202018\%20cetak\%20revisi\%20FINAL\%203juli2018.pdf}
\end{itemize}
\end{frame}

\begin{frame}[plain]
\centering
{\LARGE Thank You}\\[1.5em]
{\large Rudi Adha Prihandoko}\\[0.5em]
Department of Mathematics, Universitas Gadjah Mada
\end{frame}
\end{document}
